% !TeX program = xelatex
\documentclass[12pt,a4paper]{report}

% ============================================================
% GÓI LỆNH
% ============================================================
\usepackage{iftex}
\ifPDFTeX
  \usepackage[utf8]{inputenc}
  \usepackage[T5]{fontenc}
  \usepackage[english,vietnamese]{babel}
\else
  \usepackage{fontspec}
  \IfFontExistsTF{DejaVu Serif}
    {\setmainfont{DejaVu Serif}}
    {\setmainfont{Times New Roman}}
  \IfFontExistsTF{DejaVu Sans}
    {\setsansfont{DejaVu Sans}}
    {\setsansfont{Arial}}
  \IfFontExistsTF{DejaVu Sans Mono}
    {\setmonofont{DejaVu Sans Mono}}
    {\setmonofont{Courier New}}
  \usepackage{polyglossia}
  \setmainlanguage{vietnamese}
  \setotherlanguage{english}
\fi
\usepackage{amsmath, amssymb, amsthm}
\usepackage{mathtools}
\usepackage{geometry}
\usepackage{graphicx}
\usepackage{xcolor}
\usepackage{hyperref}
\usepackage{booktabs}
\usepackage{array}
\usepackage{enumitem}
\usepackage{fancyhdr}
\usepackage{titlesec}
\usepackage{tcolorbox}
\usepackage{listings}
\usepackage{algorithm}
\usepackage{algpseudocode}
\usepackage{mdframed}
\usepackage{caption}
\usepackage{subcaption}
\usepackage{bm}
\usepackage{tikz}
\usepackage{pgfplots}
\pgfplotsset{compat=1.18}

% ============================================================
% KÍCH THƯỚC TRANG
% ============================================================
\geometry{
  a4paper,
  top=2.5cm,
  bottom=2.5cm,
  left=3cm,
  right=2.5cm,
  headheight=15pt
}

% ============================================================
% MÀU SẮC
% ============================================================
\definecolor{mainblue}{RGB}{0,70,140}
\definecolor{lightblue}{RGB}{220,235,255}
\definecolor{darkgray}{RGB}{50,50,50}
\definecolor{accentred}{RGB}{180,0,30}
\definecolor{codegreen}{RGB}{0,120,0}
\definecolor{codebg}{RGB}{248,248,248}
\definecolor{codepurple}{RGB}{100,0,150}

% ============================================================
% SIÊU LIÊN KẾT
% ============================================================
\hypersetup{
  colorlinks=true,
  linkcolor=mainblue,
  citecolor=accentred,
  urlcolor=mainblue,
  pdfauthor={Báo cáo tổng hợp},
  pdftitle={Hàm của Ma Trận -- Báo cáo chuyên sâu},
}

% ============================================================
% ĐẦU TRANG / CHÂN TRANG
% ============================================================
\pagestyle{fancy}
\fancyhf{}
\fancyhead[L]{\small\textcolor{mainblue}{\textit{Hàm của Ma Trận (Functions of Matrices)}}}
\fancyhead[R]{\small\textcolor{mainblue}{\thepage}}
\fancyfoot[C]{\small\textcolor{gray}{Golub \& Van Loan, \textit{Matrix Computations} -- Chương 9}}
\renewcommand{\headrulewidth}{0.4pt}
\renewcommand{\footrulewidth}{0.4pt}

% ============================================================
% ĐỊNH DẠNG CHƯƠNG VÀ MỤC
% ============================================================
\titleformat{\chapter}[display]
  {\Large\bfseries\color{mainblue}}
  {\large\color{mainblue!70!black}\chaptername\ \thechapter}
  {10pt}
  {\Huge\bfseries\color{mainblue}}
  [\vspace{4pt}\titlerule\vspace{6pt}]

\titlespacing*{\chapter}{0pt}{10pt}{30pt}

\titleformat{\section}{\Large\bfseries\color{mainblue}}{{\thesection}}{1em}{}[\titlerule]
\titleformat{\subsection}{\large\bfseries\color{mainblue!80!black}}{{\thesubsection}}{1em}{}
\titleformat{\subsubsection}{\normalsize\bfseries\color{darkgray}}{{\thesubsubsection}}{1em}{}

% ============================================================
% MÔI TRƯỜNG ĐỊNH LÝ
% ============================================================
\tcbuselibrary{theorems, skins, breakable}

\newtcbtheorem[number within=section]{theorem}{Theorem}{
  enhanced, breakable,
  lines before break=0,
  colback=lightblue,
  colframe=mainblue,
  fonttitle=\bfseries,
  separator sign={.},
  label type=theorem
}{thm}

\newtcbtheorem[number within=section]{lemma}{Lemma}{
  enhanced, breakable,
  lines before break=0,
  colback=blue!5,
  colframe=blue!60!black,
  fonttitle=\bfseries,
  separator sign={.}
}{lem}

\newtcbtheorem[number within=section]{corollary}{Corollary}{
  enhanced, breakable,
  lines before break=0,
  colback=green!5,
  colframe=green!50!black,
  fonttitle=\bfseries,
  separator sign={.}
}{cor}

\newtcbtheorem[number within=section]{definition}{Definition}{
  enhanced, breakable,
  lines before break=0,
  colback=orange!8,
  colframe=orange!70!black,
  fonttitle=\bfseries,
  separator sign={.}
}{def}

\newtcbtheorem[number within=section]{remark}{Remark}{
  enhanced, breakable,
  lines before break=0,
  colback=gray!10,
  colframe=gray!60,
  fonttitle=\bfseries,
  separator sign={.}
}{rem}

% Khung thuật toán
\newmdenv[
  linecolor=mainblue,
  backgroundcolor=codebg,
  linewidth=1.2pt,
  roundcorner=4pt,
  innertopmargin=8pt,
  innerbottommargin=8pt,
  innerleftmargin=10pt,
  innerrightmargin=10pt
]{algobox}

% Khung ghi chú quan trọng
\newmdenv[
  linecolor=accentred,
  backgroundcolor=red!4,
  linewidth=1.5pt,
  roundcorner=4pt,
  innertopmargin=8pt,
  innerbottommargin=8pt,
  innerleftmargin=10pt,
  innerrightmargin=10pt
]{importantbox}

% ============================================================
% LỆNH TOÁN HỌC
% ============================================================
\newcommand{\R}{\mathbb{R}}
\newcommand{\C}{\mathbb{C}}
\newcommand{\N}{\mathbb{N}}
\newcommand{\norm}[1]{\left\lVert #1 \right\rVert}
\newcommand{\abs}[1]{\left\lvert #1 \right\rvert}
\newcommand{\diag}{\operatorname{diag}}
\newcommand{\ran}{\operatorname{ran}}
\newcommand{\sign}{\operatorname{sign}}
\newcommand{\spec}{\lambda}
\newcommand{\eA}{e^{A}}
\newcommand{\fl}{\operatorname{fl}}
\DeclareMathOperator{\re}{Re}
\DeclareMathOperator{\im}{Im}
\DeclareMathOperator{\tr}{tr}
\DeclareMathOperator{\rank}{rank}

% ============================================================
% HIỂN THỊ MÃ NGUỒN (Python)
% ============================================================
\lstdefinestyle{python}{
  language=Python,
  basicstyle=\small\ttfamily,
  keywordstyle=\color{mainblue}\bfseries,
  commentstyle=\color{codegreen}\itshape,
  stringstyle=\color{accentred},
  numberstyle=\tiny\color{gray},
  backgroundcolor=\color{codebg},
  frame=single,
  framerule=0.4pt,
  rulecolor=\color{mainblue!40},
  numbers=left,
  numbersep=8pt,
  breaklines=true,
  captionpos=b,
  tabsize=4,
  showstringspaces=false,
  morekeywords={import, from, as, None, True, False, with, yield},
  literate=
    {á}{{á}}1 {à}{{à}}1 {ả}{{ả}}1 {ã}{{ã}}1 {ạ}{{ạ}}1
    {ă}{{ă}}1 {ắ}{{ắ}}1 {ằ}{{ằ}}1 {ẳ}{{ẳ}}1 {ẵ}{{ẵ}}1 {ặ}{{ặ}}1
    {â}{{â}}1 {ấ}{{ấ}}1 {ầ}{{ầ}}1 {ẩ}{{ẩ}}1 {ẫ}{{ẫ}}1 {ậ}{{ậ}}1
    {é}{{é}}1 {è}{{è}}1 {ẻ}{{ẻ}}1 {ẽ}{{ẽ}}1 {ẹ}{{ẹ}}1
    {ê}{{ê}}1 {ế}{{ế}}1 {ề}{{ề}}1 {ể}{{ể}}1 {ễ}{{ễ}}1 {ệ}{{ệ}}1
    {í}{{í}}1 {ì}{{ì}}1 {ỉ}{{ỉ}}1 {ĩ}{{ĩ}}1 {ị}{{ị}}1
    {ó}{{ó}}1 {ò}{{ò}}1 {ỏ}{{ỏ}}1 {õ}{{õ}}1 {ọ}{{ọ}}1
    {ô}{{ô}}1 {ố}{{ố}}1 {ồ}{{ồ}}1 {ổ}{{ổ}}1 {ỗ}{{ỗ}}1 {ộ}{{ộ}}1
    {ơ}{{ơ}}1 {ớ}{{ớ}}1 {ờ}{{ờ}}1 {ở}{{ở}}1 {ỡ}{{ỡ}}1 {ợ}{{ợ}}1
    {ú}{{ú}}1 {ù}{{ù}}1 {ủ}{{ủ}}1 {ũ}{{ũ}}1 {ụ}{{ụ}}1
    {ư}{{ư}}1 {ứ}{{ứ}}1 {ừ}{{ừ}}1 {ử}{{ử}}1 {ữ}{{ữ}}1 {ự}{{ự}}1
    {ý}{{ý}}1 {ỳ}{{ỳ}}1 {ỷ}{{ỷ}}1 {ỹ}{{ỹ}}1 {ỵ}{{ỵ}}1 {đ}{{đ}}1
    {Á}{{Á}}1 {À}{{À}}1 {Ả}{{Ả}}1 {Ã}{{Ã}}1 {Ạ}{{Ạ}}1
    {Ă}{{Ă}}1 {Ắ}{{Ắ}}1 {Ằ}{{Ằ}}1 {Ẳ}{{Ẳ}}1 {Ẵ}{{Ẵ}}1 {Ặ}{{Ặ}}1
    {Â}{{Â}}1 {Ấ}{{Ấ}}1 {Ầ}{{Ầ}}1 {Ẩ}{{Ẩ}}1 {Ẫ}{{Ẫ}}1 {Ậ}{{Ậ}}1
    {É}{{É}}1 {È}{{È}}1 {Ẻ}{{Ẻ}}1 {Ẽ}{{Ẽ}}1 {Ẹ}{{Ẹ}}1
    {Ê}{{Ê}}1 {Ế}{{Ế}}1 {Ề}{{Ề}}1 {Ể}{{Ể}}1 {Ễ}{{Ễ}}1 {Ệ}{{Ệ}}1
    {Í}{{Í}}1 {Ì}{{Ì}}1 {Ỉ}{{Ỉ}}1 {Ĩ}{{Ĩ}}1 {Ị}{{Ị}}1
    {Ó}{{Ó}}1 {Ò}{{Ò}}1 {Ỏ}{{Ỏ}}1 {Õ}{{Õ}}1 {Ọ}{{Ọ}}1
    {Ô}{{Ô}}1 {Ố}{{Ố}}1 {Ồ}{{Ồ}}1 {Ổ}{{Ổ}}1 {Ỗ}{{Ỗ}}1 {Ộ}{{Ộ}}1
    {Ơ}{{Ơ}}1 {Ớ}{{Ớ}}1 {Ờ}{{Ờ}}1 {Ở}{{Ở}}1 {Ỡ}{{Ỡ}}1 {Ợ}{{Ợ}}1
    {Ú}{{Ú}}1 {Ù}{{Ù}}1 {Ủ}{{Ủ}}1 {Ũ}{{Ũ}}1 {Ụ}{{Ụ}}1
    {Ư}{{Ư}}1 {Ứ}{{Ứ}}1 {Ừ}{{Ừ}}1 {Ử}{{Ử}}1 {Ữ}{{Ữ}}1 {Ự}{{Ự}}1
    {Ý}{{Ý}}1 {Ỳ}{{Ỳ}}1 {Ỷ}{{Ỷ}}1 {Ỹ}{{Ỹ}}1 {Ỵ}{{Ỵ}}1 {Đ}{{Đ}}1,
}

\lstset{style=python}

% ============================================================
% BẮT ĐẦU TÀI LIỆU
% ============================================================
\begin{document}

% ============================================================
%  MỤC LỤC
% ============================================================
\tableofcontents
\newpage

% ============================================================
% CHƯƠNG MỞ ĐẦU: GIỚI THIỆU TỔNG QUAN
% (Không thuộc Chương 1 hay Chương 2)
% ============================================================
\chapter*{Giới thiệu tổng quan}
\addcontentsline{toc}{chapter}{Giới thiệu tổng quan}

Tính toán hàm $f(A)$ của một ma trận vuông $A \in \R^{n \times n}$ (hoặc $\C^{n \times n}$)
là một bài toán xuất hiện phổ biến trong nhiều lĩnh vực ứng dụng: hệ thống điều khiển,
phương trình vi phân ma trận, cơ học lượng tử, học máy và xử lý tín hiệu. Đây không
chỉ là sự mở rộng hình thức của khái niệm hàm số thực lên không gian ma trận, mà còn
là một lĩnh vực nghiên cứu sâu sắc với nhiều thách thức về lý thuyết lẫn tính toán số.

\section*{Bối cảnh lịch sử và tầm quan trọng}

Lịch sử của hàm ma trận gắn liền với sự phát triển của đại số tuyến tính hiện đại.
Vào cuối thế kỷ XIX, Arthur Cayley (1858) và William Rowan Hamilton đã đặt nền móng
với định lý Cayley--Hamilton: mọi ma trận đều thỏa đa thức đặc trưng của chính nó,
tức là $p(A) = 0$ với $p(\lambda) = \det(\lambda I - A)$. Từ đó, mọi hàm của $A$ có
thể được rút gọn thành đa thức bậc không quá $n-1$, cung cấp một cơ sở lý thuyết
quan trọng cho toàn bộ lĩnh vực.

Bước tiến lớn tiếp theo đến từ lý thuyết giải tích hàm (functional calculus), được
phát triển mạnh trong nửa đầu thế kỷ XX. Frigyes Riesz và Béla Sz.-Nagy đã xây dựng
lý thuyết phổ cho toán tử tuyến tính trên không gian Hilbert, từ đó lý thuyết hàm
ma trận trở thành trường hợp hữu hạn chiều đặc biệt của một khung lý thuyết rộng
lớn hơn nhiều. Đặc biệt, tích phân Cauchy cung cấp một cách định nghĩa thống nhất
và tổng quát cho mọi loại hàm ma trận.

Về mặt ứng dụng thực tiễn, nhu cầu tính toán hàm ma trận tăng vọt từ những năm
1960--1970 khi máy tính điện tử trở nên phổ biến. Cleve Moler và Charles Van Loan
(1978) tổng hợp 19 phương pháp tính hàm mũ ma trận, cho ra bài báo nổi tiếng
\emph{``Nineteen Dubious Ways to Compute the Exponential of a Matrix''} và cho thấy
không có phương pháp nào hoàn toàn vượt trội. Công trình này đã thúc đẩy hàng thập
kỷ nghiên cứu sau đó, đặc biệt bởi Nicholas Higham và các cộng sự.

\section*{Ý tưởng cơ bản}

Nếu hàm số vô hướng $f(z)$ xác định trên phổ $\lambda(A)$ của $A$, thì $f(A)$ được
xác định bằng cách thay ``$A$'' vào ``$z$'' trong công thức của $f(z)$. Ví dụ:
\begin{itemize}
  \item Nếu $f(z) = \frac{1+z}{1-z}$ và $1 \notin \lambda(A)$, thì
        \[
          f(A) = (I+A)(I-A)^{-1}.
        \]
  \item Hàm mũ: $e^A = \sum_{k=0}^{\infty} \dfrac{A^k}{k!}$.
  \item Hàm logarit: $\log(A)$ là ma trận $X$ thỏa $e^X = A$.
  \item Căn bậc hai: $A^{1/2}$ là ma trận $X$ thỏa $X^2 = A$.
\end{itemize}

Tuy nhiên, sự chuyển đổi từ hàm vô hướng sang hàm ma trận không phải lúc nào cũng
đơn giản. Khi $A$ là ma trận chéo hóa được, ta chỉ cần áp dụng $f$ lên từng trị
riêng. Nhưng khi $A$ có trị riêng lặp hoặc không chéo hóa được, định nghĩa cần
được mở rộng thông qua đạo hàm của $f$, dẫn đến lý thuyết về dạng chính tắc Jordan.

Bài toán trở nên thú vị và thách thức hơn khi $f$ là hàm siêu việt (transcendental)
như hàm mũ, logarit, lượng giác. Trong những trường hợp này, không tồn tại công
thức đóng đơn giản; ta cần dùng đến các phương pháp số với phân tích sai số cẩn thận.

\section*{Các lĩnh vực ứng dụng}

Sự quan trọng của hàm ma trận thể hiện rõ ở tính phổ biến của chúng trong nhiều
ngành khoa học và kỹ thuật:

\begin{description}[leftmargin=2cm, labelindent=0.5cm]
  \item[\textbf{Phương trình vi phân}]
    Nghiệm của hệ phương trình vi phân tuyến tính $\dot{x}(t) = Ax(t)$, $x(0) = x_0$
    có dạng $x(t) = e^{At} x_0$. Hàm mũ ma trận $e^{At}$ do đó là trung tâm của lý
    thuyết hệ động lực tuyến tính.

  \item[\textbf{Lý thuyết điều khiển}]
    Ma trận kiểm soát (controllability matrix), ma trận quan sát\\ (observability matrix),
    và các bài toán tối ưu LQR đều liên quan đến hàm mũ ma trận. Phương trình Riccati
    ma trận, nền tảng của điều khiển tối ưu, được giải thông qua hàm dấu và phân tích
    cực ma trận.

  \item[\textbf{Khoa học mạng (Network Science)}]
    Số mũ ma trận được dùng để đo tính kết nối và phân tích lan truyền thông tin trong
    mạng lưới. Centrality measures như communi\-cability dựa trực tiếp trên $e^A$ với $A$
    là ma trận kề của đồ thị.

  \item[\textbf{Cơ học lượng tử}]
    Toán tử tiến hóa trong cơ học lượng tử có dạng $U(t) = e^{-iHt/\hbar}$, trong đó
    $H$ là toán tử Hamiltonian. Việc tính toán hiệu quả ma trận $e^{-iH}$ là bài toán
    trung tâm trong mô phỏng lượng tử.

  \item[\textbf{Học máy và thống kê}]
    Trong học biểu diễn trên đa tạp Lie (Lie group learning), các phép tính trên không
    gian Riemannian đòi hỏi dùng hàm mũ và logarit ma trận. Hồi quy trên không gian
    đối xứng dương định (SPD matrices), quan trọng trong xử lý ảnh y tế và học sâu,
    dựa trên metric Riemannian với $\log$ và $\exp$ ma trận.

  \item[\textbf{Tài chính định lượng}]
    Các mô hình tín dụng liên tục (continuous-time credit models) và mô hình xác suất
    chuyển tiếp (transition probability models) liên quan đến tính $e^{Qt}$ với $Q$ là
    ma trận cường độ (generator matrix) của xích Markov liên tục.

  \item[\textbf{Phương trình vi phân đại số (DAEs)}]
    Hàm dấu ma trận được dùng để phân tách phổ của bó ma trận (matrix pencils), từ
    đó tách phần chỉ số hữu hạn và vô hạn trong DAEs.

  \item[\textbf{Xử lý tín hiệu và hình ảnh}]
    Biến đổi ma trận theo hàm lũy thừa phân số $A^\alpha$ ($\alpha \in (0,1)$) xuất
    hiện trong mô hình khuếch tán phân số (fractional diffusion), được ứng dụng trong
    nén ảnh và lọc tín hiệu thích nghi.
\end{description}

\section*{Thách thức tính toán}

Mặc dù định nghĩa lý thuyết của hàm ma trận tương đối rõ ràng thông qua dạng Jordan
hay tích phân Cauchy, việc tính toán số chính xác và hiệu quả đặt ra nhiều thách thức:

\textbf{1. Khuếch đại sai số làm tròn.} Do phép nhân ma trận tích lũy sai số,
các phương pháp naive như chuỗi Taylor có thể cho kết quả sai hoàn toàn dù mỗi
bước tính toán trông có vẻ đúng (hiện tượng ``catastrophic cancellation'').

\textbf{2. Tính không đơn trị.} Hàm logarit và căn bậc hai ma trận không đơn trị:
cùng một ma trận $A$ có thể có nhiều logarit hoặc nhiều căn bậc hai. Việc chọn
``nhánh chính'' (principal branch) và đảm bảo thuật toán hội tụ về đúng nhánh đó
đòi hỏi phân tích cẩn thận.

\textbf{3. Trị riêng lặp hoặc gần nhau.} Khi $A$ có trị riêng lặp, các sai phân
chia (divided differences) trong công thức Schur-Parlett có thể mất ổn định số.
Cần dùng công thức thay thế dựa trên giới hạn (limit forms) của sai phân chia.

\textbf{4. Ma trận không chính tắc (non-normal).} Với ma trận không chính tắc,
hàm mũ $e^{At}$ có thể tăng trước khi suy giảm dù mọi trị riêng có phần thực âm
--- hiện tượng được gọi là \emph{tăng trưởng quá độ} (transient growth). Phân tích
phổ cổ điển không nắm bắt được hành vi này; cần dùng lý thuyết pseudospectra.

\textbf{5. Chi phí tính toán.} Với ma trận lớn ($n$ hàng nghìn hoặc hàng triệu),
chi phí $O(n^3)$ của các phương pháp dày (dense) trở nên không khả thi. Phương
pháp Krylov subspace và phương pháp thưa (sparse methods) trở thành lựa chọn bắt buộc.

\section*{Tổng quan bốn hướng tiếp cận}

Báo cáo này trình bày bốn hướng tiếp cận chính, mỗi hướng có ưu và nhược điểm riêng:

\begin{enumerate}[label=\roman*.]
  \item \textbf{Phương pháp trị riêng} (Phần~\ref{sec:eigen}): dựa trên phân tích
    Jordan và Schur. Phân tích Jordan cho định nghĩa nền tảng; phân tích Schur
    cho thuật toán Schur-Parlett ổn định số trong thực tế. Phương pháp này phù hợp
    khi cần tính chính xác cao với ma trận có kích thước vừa phải.

  \item \textbf{Phương pháp xấp xỉ} (Phần~\ref{sec:approx}): xấp xỉ $f$ bởi hàm
    đơn giản hơn $g$ (đa thức, hàm hữu tỉ). Bao gồm chuỗi Taylor cắt cụt, xấp xỉ
    Padé, và các kỹ thuật cầu phương. Hiệu quả và linh hoạt nhưng đòi hỏi phân tích
    sai số cẩn thận để tránh hủy số.

  \item \textbf{Hàm mũ ma trận} (Phần~\ref{sec:expm}): phương pháp scaling and
    squaring kết hợp xấp xỉ Padé là thuật toán ``vàng'' cho $e^A$, được dùng trong
    hầu hết các thư viện số hiện đại (MATLAB, SciPy).
    Phần này cũng phân tích lý thuyết nhiễu loạn và số điều kiện của hàm mũ.

  \item \textbf{Hàm dấu, căn bậc hai và logarit} (Phần~\ref{sec:signlogsqrt}):
    các hàm này liên kết mật thiết với nhau qua lặp Newton và lặp Denman-Beavers.
    Phân tích cực ma trận cũng được trình bày như một ứng dụng quan trọng.
\end{enumerate}

\section*{Ký hiệu và quy ước}

Trong bài báo cáo ta dùng các ký hiệu sau:
\begin{itemize}
  \item $\lambda(A)$: tập phổ (tập trị riêng) của $A$.
  \item $\kappa_2(X) = \norm{X}_2 \norm{X^{-1}}_2$: số điều kiện theo chuẩn 2.
  \item $\norm{\cdot}_F$: chuẩn Frobenius.
  \item $u$: machine epsilon (độ chính xác máy tính), với IEEE 754 double precision
    $u \approx 2.22 \times 10^{-16}$.
  \item JCF: Jordan Canonical Form (dạng chính tắc Jordan).
  \item $Q^H$: liên hợp chuyển vị (Hermitian transpose) của $Q$.
  \item $\alpha(A) = \max\{\re(\lambda): \lambda \in \lambda(A)\}$: hoành độ phổ
    (spectral abscissa).
  \item $\sigma_i(A)$: trị kỳ dị (singular value) thứ $i$ của $A$ theo thứ tự giảm dần.
  \item $\rho(A) = \max\{|\lambda|: \lambda \in \lambda(A)\}$: bán kính phổ.
\end{itemize}


% ============================================================
% CHƯƠNG 1: LÝ THUYẾT VỀ HÀM CỦA MA TRẬN
% ============================================================
\chapter{Lý thuyết về hàm của ma trận}

% ============================================================
% MỤC 1: PHƯƠNG PHÁP TRỊ RIÊNG
% ============================================================
\section{Phương pháp trị riêng (Eigenvalue Methods)}
\label{sec:eigen}

\subsection{Định nghĩa hàm ma trận qua dạng Jordan}

\subsubsection{Dạng chính tắc Jordan}

Giả sử $A \in \C^{n \times n}$ có phân tích Jordan
\begin{equation}
  A = X \cdot \diag(J_1, \ldots, J_q) \cdot X^{-1},
  \label{eq:jcf}
\end{equation}
trong đó mỗi khối Jordan có dạng
\[
  J_i = \begin{pmatrix}
    \lambda_i & 1      & \cdots & 0      \\
    0         & \lambda_i & \ddots & \vdots \\
    \vdots    & \ddots & \ddots & 1      \\
    0         & \cdots & 0      & \lambda_i
  \end{pmatrix} \in \C^{n_i \times n_i}, \quad i = 1, \ldots, q.
\]

\begin{definition}{Hàm ma trận (Jordan)}{matfun}
Hàm $f(A)$ được xác định bởi
\[
  f(A) = X \cdot \diag(F_1, \ldots, F_q) \cdot X^{-1},
\]
trong đó
\[
  F_i = \begin{pmatrix}
    f(\lambda_i) & f^{(1)}(\lambda_i) & \cdots & \dfrac{f^{(n_i-1)}(\lambda_i)}{(n_i-1)!} \\[6pt]
    0            & f(\lambda_i)        & \ddots & \vdots                                      \\
    \vdots       & \ddots              & \ddots & f^{(1)}(\lambda_i)                          \\
    0            & \cdots              & 0      & f(\lambda_i)
  \end{pmatrix}.
\]
Định nghĩa này đòi hỏi tất cả các đạo hàm $f^{(k)}(\lambda_i)$ cần thiết phải tồn tại.
\end{definition}

\subsubsection{Biểu diễn chuỗi Taylor}

Nếu $f$ có thể biểu diễn bằng chuỗi Taylor trên phổ của $A$, thì $f(A)$ có cùng
biểu diễn đó nhưng thay $z$ bởi $A$.

\begin{theorem}{Biểu diễn chuỗi Taylor}{taylor}
Nếu $f$ có biểu diễn Taylor
\[
  f(z) = \sum_{k=0}^{\infty} \frac{f^{(k)}(z_0)}{k!}(z - z_0)^k, \quad |z - z_0| < r,
\]
và $|\lambda - z_0| < r$ với mọi $\lambda \in \lambda(A)$, thì
\[
  f(A) = \sum_{k=0}^{\infty} \frac{f^{(k)}(z_0)}{k!}(A - z_0 I)^k.
\]
\end{theorem}

Bổ đề nền tảng cho định lý này xét trường hợp $A$ là một khối Jordan đơn:

\begin{lemma}{Khối Jordan đơn}{jordanblock}
Giả sử $B \in \C^{m \times m}$ là khối Jordan, $B = \lambda I_m + E$ ($E$ là phần
bidiagonal trên nghiêm ngặt). Nếu $|\lambda - z_0| < r$, thì
\[
  f(B) = \sum_{k=0}^{\infty} \frac{f^{(k)}(z_0)}{k!}(B - z_0 I_m)^k.
\]
Hơn nữa, nếu $0 \le p \le m-1$, thì $[E^p]_{ij} = \delta_{i,j-p}$, và
\[
  f(B) = \sum_{p=0}^{m-1} \frac{f^{(p)}(\lambda)}{p!} E^p.
\]
\end{lemma}

Các hàm ma trận quan trọng có biểu diễn chuỗi Taylor đơn giản bao gồm:
\[
  e^A = \sum_{k=0}^{\infty} \frac{A^k}{k!}, \quad
  \log(I-A) = -\sum_{k=1}^{\infty} \frac{A^k}{k}\ (|\lambda| < 1),
\]
\[
  \sin(A) = \sum_{k=0}^{\infty} \frac{(-1)^k A^{2k+1}}{(2k+1)!}, \quad
  \cos(A) = \sum_{k=0}^{\infty} \frac{(-1)^k A^{2k}}{(2k)!}.
\]

Từ định nghĩa này ta có hai tính chất giao hoán quan trọng:
\[
  A \cdot f(A) = f(A) \cdot A, \qquad
  f(X^{-1}AX) = X^{-1} f(A) X.
\]

\subsection{Phương pháp vector riêng (Eigenvector Approach)}

Khi $A$ là ma trận chéo hóa được: $A = X \diag(\lambda_1, \ldots, \lambda_n) X^{-1}$,
ta có hệ quả sau:

\begin{corollary}{Ma trận chéo hóa được}{diag}
\[
  f(A) = X \cdot \diag\bigl(f(\lambda_1), \ldots, f(\lambda_n)\bigr) \cdot X^{-1}.
\]
\end{corollary}


\subsubsection{Ví dụ minh họa}

Xét
\[
  A = \begin{pmatrix} 1 + 10^{-5} & 1 \\ 0 & 1 - 10^{-5} \end{pmatrix}.
\]
Ma trận vector riêng có $\kappa_2(X) \approx 10^5$. Với $u \approx 10^{-7}$, tính
$e^A$ qua đường chéo hóa cho kết quả sai ở chữ số thứ 2 của phần tử $(1,2)$, trong
khi kết quả chính xác là $[e^A]_{12} \approx 2.718282$.

\subsection{Phương pháp phân tích Schur}

\subsubsection{Ý tưởng chung}

Một số khó khăn của phương pháp Jordan có thể được khắc phục bằng cách dùng
\emph{phân tích Schur}: $A = Q T Q^H$ ($Q$ unitary, $T$ tam giác trên). Khi đó:
\[
  f(A) = Q \, f(T) \, Q^H.
\]
Vấn đề còn lại là tính $f(T)$ cho ma trận tam giác trên.

\subsubsection{Công thức tường minh cho ma trận tam giác trên}

\begin{theorem}{Hàm của ma trận tam giác trên -- Schur}{schurfun}
Gọi $T = (t_{ij})$ là ma trận tam giác trên $n \times n$ với $\lambda_i = t_{ii}$,
$F = f(T) = (f_{ij})$. Thì $f_{ij} = 0$ khi $i > j$, $f_{ii} = f(\lambda_i)$, và với
$i < j$:
\[
  f_{ij} = \sum_{(s_0,\ldots,s_k) \in S_{ij}}
           t_{s_0 s_1} t_{s_1 s_2} \cdots t_{s_{k-1} s_k}
           \, f[\lambda_{s_0}, \ldots, \lambda_{s_k}],
\]
trong đó $S_{ij}$ là tập tất cả dãy tăng nghiêm ngặt từ $i$ đến $j$, và
$f[\lambda_{s_0},\ldots,\lambda_{s_k}]$ là \emph{sai phân chia bậc $k$} (divided difference)
của $f$.
\end{theorem}

Ví dụ cho $3\times 3$:
\[
  f\!\begin{pmatrix}\lambda_1 & t_{12} & t_{13} \\ 0 & \lambda_2 & t_{23} \\ 0 & 0 & \lambda_3\end{pmatrix}
  =
  \begin{pmatrix}
    f(\lambda_1) & t_{12}\dfrac{f(\lambda_2)-f(\lambda_1)}{\lambda_2-\lambda_1} & F_{13} \\[8pt]
    0 & f(\lambda_2) & t_{23}\dfrac{f(\lambda_3)-f(\lambda_2)}{\lambda_3-\lambda_2} \\[8pt]
    0 & 0 & f(\lambda_3)
  \end{pmatrix},
\]
trong đó $F_{13}$ là tổ hợp tuyến tính của các sai phân chia bậc 1 và bậc 2.

\subsubsection{Thuật toán Schur-Parlett}

Parlett (1974) phát triển một phương pháp đệ quy hiệu quả dựa trên phương trình
$FT = TF$. So sánh phần tử $(i,j)$ của hai vế (với $t_{ii} \ne t_{jj}$):
\[
  f_{ij} = \frac{t_{ij}(f_{jj} - f_{ii})}{t_{jj} - t_{ii}}
           + \sum_{k=i+1}^{j-1} \frac{t_{ik} f_{kj} - f_{ik} t_{kj}}{t_{jj} - t_{ii}}.
\]

\begin{algobox}
\textbf{Thuật toán 1.1.1 (Schur-Parlett)} \quad Tính $F = f(T)$ với $T$ tam giác trên,
trị riêng phân biệt.

\medskip
\textbf{Bước 1.} Khởi tạo đường chéo: $f_{ii} = f(t_{ii})$, $i = 1:n$.

\textbf{Bước 2.} Với $p = 1, 2, \ldots, n-1$ (siêu đường chéo thứ $p$):
\[
  \text{Với } i = 1:n-p,\quad j = i+p:
  \quad
  f_{ij} = \frac{t_{ij}(f_{jj}-f_{ii}) + \sum_{k=i+1}^{j-1}(t_{ik}f_{kj} - f_{ik}t_{kj})}{t_{jj} - t_{ii}}.
\]

\medskip
\textit{Độ phức tạp:} $\tfrac{2}{3}n^3$ phép tính dấu phẩy động (flops).
Khi đó $f(A) = Q F Q^H$.
\end{algobox}

\subsubsection{Phương pháp Schur-Parlett theo khối}

Khi $A$ có trị riêng lặp hoặc gần nhau, sai phân chia trong Thuật toán 1.1.1 trở
nên bất ổn định số. Giải pháp là dùng phiên bản \emph{theo khối}:

\begin{enumerate}
  \item Chọn $Q$ sao cho $T$ được phân hoạch thành các khối chéo
        $T_{11}, T_{22}, \ldots, T_{pp}$ với $\lambda(T_{ii}) \cap \lambda(T_{jj}) = \emptyset$
        và mỗi $T_{ii}$ gắn với một cụm trị riêng.
  \item Tính $F_{ii} = f(T_{ii})$ cho từng khối chéo (dùng phương pháp riêng cho
        cụm trị riêng).
  \item Tính các khối ngoài đường chéo $F_{ij}$ ($i < j$) qua hệ Sylvester:
        \[
          F_{ij}T_{jj} - T_{ii}F_{ij} = T_{ij}F_{jj} - F_{ii}T_{ij}
          + \sum_{k=i+1}^{j-1}(T_{ik}F_{kj} - F_{ik}T_{kj}).
        \]
        Hệ này được giải bằng thuật toán Bartels-Stewart.
\end{enumerate}

\subsection{Độ nhạy của hàm ma trận}

\begin{definition}{Số điều kiện tương đối}{condrel}
Số điều kiện tương đối của $f$ tại $A \in \C^{n \times n}$ là
\[
  \mathrm{cond}_{\mathrm{rel}}(f, A)
  = \lim_{\varepsilon \to 0}
    \sup_{\norm{E} \le \varepsilon \norm{A}}
    \frac{\norm{f(A+E) - f(A)}}{\norm{f(A)}}.
\]
\end{definition}

Đây thực chất là chuẩn hóa của đạo hàm Fréchet của ánh xạ $A \mapsto f(A)$.
Việc cài đặt cẩn thận thuật toán Schur-Parlett theo khối thường cho kết quả
\emph{ổn định tiến} (forward stable):
\[
  \frac{\norm{\hat{F} - f(A)}}{\norm{f(A)}} \approx u \cdot \mathrm{cond}_{\mathrm{rel}}(f,A).
\]
Điều này \emph{không thể} đảm bảo với phương pháp chéo hóa khi $X$ kém điều kiện.

% ------------------------------------------------------------------
% PYTHON MINH HỌA – PHƯƠNG PHÁP TRỊ RIÊNG
% ------------------------------------------------------------------
\subsection{Mã nguồn Python minh họa}

Đoạn mã dưới đây minh họa ba cách tính $e^A$: qua phân tích trị riêng (vector riêng),
qua phân tích Schur (Schur-Parlett đơn giản hóa), và dùng thư viện SciPy. Kết quả
so sánh cho thấy ảnh hưởng của số điều kiện $\kappa_2(X)$ lên độ chính xác.

\begin{lstlisting}[style=python, caption={Phương pháp trị riêng: tính $f(A) = e^A$ và so sánh độ chính xác}]
import numpy as np
from scipy.linalg import expm, schur

# -------------------------------------------------------
# 1. Ma trận ví dụ: có số điều kiện cao
# -------------------------------------------------------
eps = 1e-5
A = np.array([[1 + eps, 1.0],
              [0.0,     1 - eps]], dtype=float)

# -------------------------------------------------------
# 2. Phương pháp 1: Phân tích trị riêng (eigenvector)
# -------------------------------------------------------
eigenvalues, X = np.linalg.eig(A)
kappa = np.linalg.cond(X)               # số điều kiện của X
f_eig = np.exp(eigenvalues)             # áp dụng f = exp lên trị riêng
fA_eig = X @ np.diag(f_eig) @ np.linalg.inv(X)

# -------------------------------------------------------
# 3. Phương pháp 2: Phân tích Schur + công thức phần tử
# -------------------------------------------------------
T, Q = schur(A, output='complex')       # A = Q T Q^H
# Tính f(T) cho ma trận tam giác trên 2x2
fT = np.zeros_like(T)
lam = np.diag(T)
fT[0, 0] = np.exp(lam[0])
fT[1, 1] = np.exp(lam[1])
# Sai phân chia bậc 1 cho phần tử (0,1)
if abs(lam[1] - lam[0]) > 1e-14:
    fT[0, 1] = T[0, 1] * (fT[1, 1] - fT[0, 0]) / (lam[1] - lam[0])
else:
    fT[0, 1] = T[0, 1] * np.exp(lam[0])   # giới hạn khi lambda trùng nhau
fA_schur = Q @ fT @ Q.conj().T

# -------------------------------------------------------
# 4. Phương pháp 3: SciPy expm (Padé + thu nhỏ và bình phương)
# -------------------------------------------------------
fA_scipy = expm(A)

# -------------------------------------------------------
# 5. So sánh kết quả
# -------------------------------------------------------
print("="*55)
print(f"  Số điều kiện của ma trận vector riêng: kappa = {kappa:.2e}")
print("="*55)
print("\nKết quả qua phân tích trị riêng (eigenvector):")
print(fA_eig.real)
print("\nKết quả qua Schur-Parlett (đơn giản hóa):")
print(fA_schur.real)
print("\nKết quả SciPy expm (chuẩn xác):")
print(fA_scipy)

err_eig    = np.linalg.norm(fA_eig.real - fA_scipy, 'fro')
err_schur  = np.linalg.norm(fA_schur.real - fA_scipy, 'fro')
print(f"\nSai số Frobenius (eigenvector vs scipy): {err_eig:.2e}")
print(f"Sai số Frobenius (Schur      vs scipy): {err_schur:.2e}")
\end{lstlisting}

% ============================================================
% MỤC 2: PHƯƠNG PHÁP XẤP XỈ
% ============================================================
\section{Phương pháp xấp xỉ (Approximation Methods)}
\label{sec:approx}

Ý tưởng cốt lõi: nếu $g(z) \approx f(z)$ trên $\lambda(A)$, thì $g(A) \approx f(A)$.
Ví dụ:
\[
  e^A \approx I + A + \frac{A^2}{2!} + \cdots + \frac{A^q}{q!}.
\]

\subsection{Chặn sai số: Phân tích Jordan}

\begin{theorem}{Sai số xấp xỉ -- Jordan}{errjordan}
Giả sử $A = X \diag(J_1,\ldots,J_q) X^{-1}$ là JCF của $A$. Nếu $f,g$ giải tích
trên tập mở chứa $\lambda(A)$, thì
\[
  \norm{f(A) - g(A)}_2
  \le \kappa_2(X) \cdot \max_{\substack{1 \le i \le q \\ 0 \le r \le n_i-1}}
      \frac{\abs{f^{(r)}(\lambda_i) - g^{(r)}(\lambda_i)}}{r!}.
\]
Chặn này phụ thuộc trực tiếp vào số điều kiện $\kappa_2(X)$ của ma trận vector riêng.
\end{theorem}

\subsection{Chặn sai số: Phân tích Schur}

\begin{theorem}{Sai số xấp xỉ -- Schur}{errschur}
Gọi $Q^H AQ = T = \diag(\lambda_i) + N$ là phân tích Schur của $A$, $N$ là phần
tam giác trên nghiêm ngặt. Nếu $f,g$ giải tích trên tập lồi đóng $\Omega$ chứa
$\lambda(A)$ trong nội bộ, thì
\[
  \norm{f(A) - g(A)}_F
  \le \sum_{r=0}^{n-1} \delta_r \,
      \frac{\bigl\| |N|^r \bigr\|_F}{r!},
  \qquad
  \delta_r = \sup_{z \in \Omega} \abs{f^{(r)}(z) - g^{(r)}(z)}.
\]
\end{theorem}

\begin{remark}{So sánh hai chặn}{compare_bounds}
Chặn Jordan phụ thuộc $\kappa_2(X)$; chặn Schur phụ thuộc chuẩn của $N$ (độ lệch
chuẩn tắc). Với ma trận kém điều kiện eigensystem, chặn Jordan có thể quá bi quan.
Ví dụ: nếu $\kappa_2(X) \approx 10^7$ thì chặn Jordan dự đoán sai số $O(1)$ trong
khi chặn Schur cho $O(10^{-2})$ và sai số thực tế $\approx 10^{-5}$.
\end{remark}

\subsection{Xấp xỉ bằng chuỗi Taylor cắt cụt}

\begin{theorem}{Sai số Taylor cắt cụt}{taylorerr}
Nếu $f(z) = \sum_{k=0}^\infty \alpha_k z^k$ trên đĩa mở chứa $\lambda(A)$, thì
\[
  \left\| f(A) - \sum_{k=0}^{q} \alpha_k A^k \right\|_2
  \le \frac{n}{(q+1)!} \max_{0 \le s \le 1}
      \norm{A^{q+1} f^{(q+1)}(sA)}_2.
\]
\end{theorem}

\begin{importantbox}
\textbf{Cảnh báo thực tế:} Việc lấy nhiều số hạng hơn không nhất thiết tăng độ chính
xác! Các tổng trung gian có thể có phần tử cỡ lớn, dẫn đến hủy số nghiêm trọng
(catastrophic cancellation). Ví dụ với
$A = \begin{pmatrix} -49 & 24 \\ -64 & 31 \end{pmatrix}$:
lấy $q=59$ cho lỗi lý thuyết $<10^{-60}$, nhưng kết quả tính thực tế sai hoàn toàn
vì các tổng trung gian có phần tử $\sim 10^7$ vượt ngưỡng máy tính $u \approx 10^{-7}$.
\end{importantbox}

\subsubsection{Kỹ thuật thu nhỏ (scaling)}

Để khắc phục hạn chế của Taylor cắt cụt gần gốc tọa độ, người ta dùng các công thức
nhân đôi góc. Ví dụ cho sin/cos:
\[
  \cos(2A) = 2\cos^2(A) - I, \qquad \sin(2A) = 2\sin(A)\cos(A).
\]
Từ đó xây dựng $\sin(A)$ và $\cos(A)$ từ các xấp xỉ Taylor cho $\sin(A/2^k)$ và
$\cos(A/2^k)$ với $k$ đủ lớn để $\norm{A}_\infty \approx 2^k$.

\subsection{Tính đa thức ma trận hiệu quả}

Cần tính $p(A) = b_0 I + b_1 A + \cdots + b_q A^q$. Horner thông thường cần $q-1$
phép nhân ma trận. Tuy nhiên, thuật toán Paterson--Stockmeyer cho phép tính với
$s + r - 1$ phép nhân khi chọn $s = \lfloor\sqrt{q}\rfloor$, $r = \lfloor q/s \rfloor$:
\[
  p(A) = \sum_{k=0}^{r} B_k (A^s)^k,
\]
giảm đáng kể số phép nhân ma trận.

\begin{algobox}
\textbf{Thuật toán 1.2.1 (Horner scheme cho đa thức ma trận)}

\medskip
Nhập: $A \in \R^{n\times n}$, hệ số $b_0,\ldots,b_q$.
\[
  F \leftarrow b_q A + b_{q-1}I; \quad
  \text{Với } k = q-2, q-3, \ldots, 0: F \leftarrow AF + b_k I.
\]
Độ phức tạp: $q-1$ phép nhân ma trận.
\end{algobox}

\subsection{Lũy thừa ma trận nhị phân (Binary Powering)}

\begin{algobox}
\textbf{Thuật toán 1.2.2 (Binary Powering)} \quad Tính $F = A^s$.

\medskip
Viết $s = \sum_{k=0}^{t} \beta_k 2^k$ (khai triển nhị phân). Bắt đầu từ $Z = A$,
bình phương liên tiếp và nhân có chọn lọc theo bit 1.

\medskip
Độ phức tạp: tối đa $2\lfloor \log_2 s \rfloor$ phép nhân ma trận.
\end{algobox}

\subsection{Tích phân hàm ma trận}

Để tính $F = \int_a^b f(At)\,dt$, áp dụng quy tắc cầu phương (ví dụ: Simpson):
\[
  \tilde{F} = \frac{h}{3}\sum_{k=0}^{m} w_k f(A(a + kh)),
\]
với sai số bị chặn bởi
\[
  \norm{\tilde{F} - F}_2 \le \frac{n h^4 (b-a)}{180}
  \max_{a \le t \le b} \norm{f^{(4)}(At)}_2.
\]

\subsection{Biểu diễn tích phân Cauchy}

Còn một cách định nghĩa $f(A)$ khác thông qua định lý tích phân Cauchy:
\[
  f(A) = \frac{1}{2\pi i} \oint_\Gamma f(z)(zI - A)^{-1}\,dz,
\]
trong đó $\Gamma$ là đường cong kín bao quanh $\lambda(A)$ và $f$ giải tích trong
và trên $\Gamma$. Phương pháp này, kết hợp với công cụ cầu phương bậc cao, cho phép
tính toán một số loại hàm ma trận hiệu quả.

% ------------------------------------------------------------------
% PYTHON MINH HỌA – PHƯƠNG PHÁP XẤP XỈ
% ------------------------------------------------------------------
\subsection{Mã nguồn Python minh họa}

Mã nguồn dưới đây trình bày: (1) xấp xỉ Taylor cắt cụt và hiện tượng hủy số,
(2) lược đồ Horner tính đa thức ma trận, (3) so sánh sai số theo bậc xấp xỉ $q$.

\begin{lstlisting}[style=python, caption={Phương pháp xấp xỉ: Taylor cắt cụt, Horner, và hiện tượng catastrophic cancellation}]
import numpy as np
from scipy.linalg import expm

# -------------------------------------------------------
# Hàm tính e^A bằng chuỗi Taylor cắt cụt bậc q
# -------------------------------------------------------
def expm_taylor(A, q):
    """Xấp xỉ e^A bằng tổng Taylor bậc q."""
    n = A.shape[0]
    F    = np.eye(n)          # số hạng k=0
    Ak   = np.eye(n)          # A^k
    fact = 1.0
    for k in range(1, q + 1):
        Ak   = Ak @ A
        fact *= k
        F   += Ak / fact
    return F

# Hàm tính đa thức ma trận p(A) bằng lược đồ Horner
def horner(A, coeffs):
    """
    Tính p(A) = coeffs[0]*I + coeffs[1]*A + ... + coeffs[q]*A^q
    dùng lược đồ Horner.
    coeffs[0] là hệ số tự do, coeffs[-1] là hệ số bậc cao nhất.
    """
    n = A.shape[0]
    q = len(coeffs) - 1
    F = coeffs[q] * np.eye(n)
    for k in range(q - 1, -1, -1):
        F = F @ A + coeffs[k] * np.eye(n)
    return F

# -------------------------------------------------------
# Ví dụ 1: Ma trận hiền lành
# A có trị riêng âm (ổn định, hàm mũ giảm)
# -------------------------------------------------------
A_ok = np.array([[1.0, 2.0],
                 [0.0, 3.0]])
ref  = expm(A_ok)

print("Ví dụ 1 -- Ma trận bình thường A = [[1,2],[0,3]]")
print(f"  Kết quả chuẩn xác (SciPy): \n{ref}\n")
for q in [5, 10, 20, 30]:
    approx = expm_taylor(A_ok, q)
    err = np.linalg.norm(approx - ref, 'fro')
    print(f"  q={q:2d}  ||Taylor - chính xác||_F = {err:.2e}")

# -------------------------------------------------------
# Ví dụ 2: Ma trận "ngổ ngáo" -- hủy số nghiêm trọng
# A có trị riêng 1 và -17 (phổ lớn)
# -------------------------------------------------------
A_bad = np.array([[-49.0, 24.0],
                  [-64.0, 31.0]])
ref_bad = expm(A_bad)

print("\nVí dụ 2 -- Ma trận gây hủy số A = [[-49,24],[-64,31]]")
print(f"  Kết quả chuẩn xác:\n{ref_bad}\n")
print("  Bậc q  ||Taylor - chính xác||_F   Max|phần tử trung gian|")
for q in [10, 20, 40, 59]:
    approx = expm_taylor(A_bad, q)
    err = np.linalg.norm(approx - ref_bad, 'fro')
    # Tính giá trị lớn nhất của tổng trung gian (theo dõi hủy số)
    n = A_bad.shape[0]
    Ak = np.eye(n); fact = 1.0; max_elem = 0.0
    for k in range(1, q + 1):
        Ak = Ak @ A_bad; fact *= k
        max_elem = max(max_elem, np.max(np.abs(Ak / fact)))
    print(f"  q={q:3d}  {err:.2e}                  {max_elem:.2e}")

# -------------------------------------------------------
# Ví dụ 3: Lược đồ Horner cho đa thức ma trận
# p(x) = 1 + x + x^2/2 + x^3/6 (Taylor bậc 3 của e^x)
# -------------------------------------------------------
A_small = np.array([[0.1, 0.2],
                    [0.0, 0.3]])
coeffs = [1.0, 1.0, 0.5, 1/6]   # hệ số [b0, b1, b2, b3]
p_horner = horner(A_small, coeffs)
p_direct = (np.eye(2) + A_small
            + A_small @ A_small / 2
            + A_small @ A_small @ A_small / 6)
print("\nVí dụ 3 -- Kiểm tra Horner vs tính trực tiếp:")
print(f"  ||Horner - tính trực tiếp||_F = {np.linalg.norm(p_horner - p_direct):.2e}")
\end{lstlisting}

% ============================================================
% MỤC 3: HÀM MŨ MA TRẬN
% ============================================================
\section{Hàm mũ ma trận (The Matrix Exponential)}
\label{sec:expm}

Hàm mũ ma trận
\[
  e^{At} = \sum_{k=0}^{\infty} \frac{(At)^k}{k!}
\]
là một trong những hàm ma trận được tính toán nhiều nhất trong thực tế, đặc biệt
trong lý thuyết điều khiển và giải phương trình vi phân.

\subsection{Phương pháp xấp xỉ Padé kết hợp Scaling and Squaring}

\subsubsection{Hàm Padé}

Xấp xỉ Padé $(p,q)$ cho hàm mũ được định nghĩa bởi:
\[
  R_{pq}(z) = D_{pq}(z)^{-1} N_{pq}(z),
\]
trong đó
\[
  N_{pq}(z) = \sum_{k=0}^{p} \frac{(p+q-k)!\,p!}{(p+q)!\,k!\,(p-k)!} z^k,
  \qquad
  D_{pq}(z) = \sum_{k=0}^{q} \frac{(p+q-k)!\,q!}{(p+q)!\,k!\,(q-k)!} (-z)^k.
\]
Trường hợp $q=0$: $R_{p0}(z) = 1 + z + \cdots + z^p/p!$ (đa thức Taylor bậc $p$).

\subsubsection{Hạn chế và cách khắc phục}

Xấp xỉ Padé chỉ chính xác gần gốc. Khai triển sai số:
\[
  e^A = R_{pq}(A)
  + \frac{(-1)^q}{(p+q)!} A^{p+q+1} D_{pq}(A)^{-1}
    \int_0^1 u^p(1-u)^q e^{A(1-u)}\,du.
\]
Khắc phục bằng kỹ thuật \emph{thu nhỏ và bình phương} (scaling and squaring):
\[
  e^A = \left(e^{A/m}\right)^m, \quad m = 2^j \text{ (chọn $j$ phù hợp)}.
\]
Tính $F_{pq} = R_{pq}(A/2^j) \approx e^{A/2^j}$, rồi bình phương $j$ lần:
$e^A \approx F_{pq}^{2^j}$.

\subsubsection{Chặn sai số}

Nếu $\norm{A}_\infty / 2^j \le 1/2$, tồn tại $E$ sao cho $F_{pq} = e^{A+E}$, $AE=EA$, và
\[
  \norm{E}_\infty \le \varepsilon(p,q)\norm{A}_\infty,
  \quad
  \varepsilon(p,q) = 2^{3-(p+q)} \frac{p!\,q!}{(p+q)!\,(p+q+1)!}.
\]

\begin{algobox}
\textbf{Thuật toán 1.3.1 (Scaling and Squaring)} \\
Nhập: $\delta > 0$, $A \in \R^{n\times n}$. Xuất: $F \approx e^A$ với sai số $\le \delta \norm{A}_\infty$.

\medskip
1. $j = \max\{0,\, 1+\lfloor \log_2\norm{A}_\infty \rfloor\}$; \quad $A \leftarrow A/2^j$.\\
2. Chọn $q$ nhỏ nhất sao cho $\varepsilon(q,q) \le \delta$.\\
3. Tính $N = N_{qq}(A)$ và $D = D_{qq}(A)$ bằng Horner.\\
4. Giải $DF = N$ bằng phép khử Gauss.\\
5. Bình phương $j$ lần: $F \leftarrow F^2$ ($j$ lần).

\medskip
\textit{Độ phức tạp:} $\approx 2(q + j + 1/3)n^3$ flops.
\end{algobox}

\subsection{Tối ưu hóa tính đa thức Padé}

Dùng kỹ thuật Horner đặc biệt của §1.2.4, ví dụ với $q=8$:
\[
  N_{88}(A) = U + AV, \quad D_{88}(A) = U - AV,
\]
với $U = c_0 I + c_2 A^2 + (c_4 I + c_6 A^2 + c_8 A^4)A^4$,
$V = c_1 I + c_3 A^2 + (c_5 I + c_7 A^2)A^4$.
Chỉ cần 5 phép nhân thay vì 7.

\subsection{Lý thuyết nhiễu loạn của hàm mũ ma trận}

Xuất phát điểm là bài toán giá trị ban đầu:
\[
  \dot{X}(t) = AX(t), \quad X(0) = I \implies X(t) = e^{At}.
\]
Đồng nhất thức quan trọng:
\[
  e^{(A+E)t} - e^{At} = \int_0^t e^{A(t-s)} E\, e^{(A+E)s}\,ds.
\]

\subsubsection{Chặn tăng trưởng qua phân tích Schur}

Gọi $Q^H AQ = \diag(\lambda_i) + N$ là phân tích Schur. Thì:
\[
  \norm{e^{At}}_2 \le e^{\alpha(A)t} M_S(t), \quad
  M_S(t) = \sum_{k=0}^{n-1} \frac{(\norm{N}_2 t)^k}{k!},
\]
trong đó $\alpha(A) = \max\{\re(\lambda): \lambda \in \lambda(A)\}$ là \emph{hoành độ phổ}
(spectral abscissa). Nếu $A$ là chính tắc ($N=0$), thì $M_S(t) \equiv 1$ và
$\norm{e^{At}}_2 = e^{\alpha(A)t}$.

\subsubsection{Số điều kiện của hàm mũ ma trận}

\begin{definition}{Số điều kiện $\nu(A,t)$}{condexp}
\[
  \nu(A,t) = \max_{\norm{E} \le 1}
  \frac{\displaystyle\left\|\int_0^t e^{A(t-s)} E\, e^{As}\,ds\right\|_2}
       {\norm{A}_2 \norm{e^{At}}_2}.
\]
\end{definition}

Tính chất: $\nu(A,t) \ge t\norm{A}_2$, với đẳng thức khi và chỉ khi $A$ chính tắc.
Nếu $\nu(A,t)$ lớn, thay đổi nhỏ trong $A$ có thể gây thay đổi lớn trong $e^{At}$.

\subsection{pseudospectra và hành vi quá độ}

Nếu $A$ không chính tắc, hàm mũ có thể \emph{tăng trước khi suy giảm} dù mọi trị riêng
đều có phần thực âm. Ví dụ minh họa:
\[
  A = \begin{pmatrix} -1 & 1000 \\ 0 & -1 \end{pmatrix}
  \implies
  e^{At} = e^{-t}\begin{pmatrix} 1 & 1000t \\ 0 & 1 \end{pmatrix}.
\]
$\norm{e^{At}}_2$ đạt đỉnh $\approx 500$ tại $t \approx 1$ dù $\alpha(A) = -1 < 0$.

pseudospectra cho chặn dưới hữu dụng: với mọi $\varepsilon > 0$,
\[
  \sup_{t > 0} \norm{e^{At}}_2 \ge \frac{\alpha_\varepsilon(A)}{\varepsilon},
\]
trong đó $\alpha_\varepsilon(A) = \sup_{z \in \Lambda_\varepsilon(A)} \re(z)$ là
\emph{hoành độ pseudospectra $\varepsilon$}.

\subsection{Phân tích ổn định của thuật toán Scaling and Squaring}

Vấn đề trong bước bình phương: nếu $\norm{e^{At}}$ tăng mạnh ban đầu, sai số làm tròn
tích lũy qua các lần bình phương:
\[
  \gamma = u \cdot \norm{G^2}_2 \cdot \norm{G^4}_2 \cdots \norm{G^{2^{j-1}}}_2.
\]
Nếu có tăng trưởng ban đầu lớn, $\gamma \gg u\norm{e^A}_2$, tức là không thể đảm bảo
sai số tương đối nhỏ. Với $A$ chính tắc, $\gamma \approx u\norm{e^A}_2$: thuật toán
luôn cho sai số tương đối nhỏ.

% ------------------------------------------------------------------
% PYTHON MINH HỌA – HÀM MŨ MA TRẬN
% ------------------------------------------------------------------
\subsection{Mã nguồn Python minh họa}

Mã dưới đây cài đặt thuật toán Scaling and Squaring với xấp xỉ Padé bậc $(6,6)$,
sau đó minh họa hiện tượng tăng trưởng quá độ với ma trận không chính tắc.

\begin{lstlisting}[style=python, caption={Hàm mũ ma trận: Scaling and Squaring + Padé, và hiện tượng transient growth}]
import numpy as np
from scipy.linalg import expm
from math import factorial

# -------------------------------------------------------
# Xấp xỉ Padé (p,p) cho hàm mũ
# -------------------------------------------------------
def pade_coeff(p):
    """Trả về hệ số của N_pp(z) và D_pp(z)."""
    c = np.zeros(p + 1)
    for k in range(p + 1):
        c[k] = (factorial(p + p - k) * factorial(p)
                / (factorial(p + p) * factorial(k) * factorial(p - k)))
    # N_pp(z) = sum c[k] z^k
    # D_pp(z) = sum c[k] (-z)^k  =>  hệ số: c[k]*(-1)^k
    N_coeffs = c.copy()
    D_coeffs = c * np.array([(-1)**k for k in range(p + 1)])
    return N_coeffs, D_coeffs

def poly_matrix(A, coeffs):
    """Tính đa thức ma trận bằng Horner: tổng theo k của coeffs[k] A^k."""
    n = A.shape[0]
    q = len(coeffs) - 1
    F = coeffs[q] * np.eye(n)
    for k in range(q - 1, -1, -1):
        F = F @ A + coeffs[k] * np.eye(n)
    return F

def expm_scaling_squaring(A, p=6):
    """
    Scaling and Squaring với xấp xỉ Padé bậc (p,p).
    - Thu nhỏ: chọn j sao cho ||A/2^j||_inf <= 0.5
    - Tính Padé: R = D^{-1} N
    - Bình phương j lần
    """
    A = np.array(A, dtype=float)
    norm_inf = np.max(np.sum(np.abs(A), axis=1))
    j = max(0, int(np.ceil(np.log2(norm_inf / 0.5)))) if norm_inf > 0.5 else 0
    As = A / (2 ** j)                      # thu nhỏ

    N_c, D_c = pade_coeff(p)
    N = poly_matrix(As, N_c)
    D = poly_matrix(As, D_c)
    F = np.linalg.solve(D, N)              # F = D^{-1} N

    for _ in range(j):                     # bình phương j lần
        F = F @ F
    return F

# -------------------------------------------------------
# So sánh với SciPy expm
# -------------------------------------------------------
A_test = np.array([[2.0, 1.0, 0.0],
                   [0.0, -1.0, 3.0],
                   [0.0, 0.0, 0.5]])
F_mine  = expm_scaling_squaring(A_test, p=6)
F_scipy = expm(A_test)
err = np.linalg.norm(F_mine - F_scipy, 'fro')
print("So sánh Scaling-Squaring (Padé p=6) vs SciPy expm:")
print(f"  ||F_mine - F_scipy||_F = {err:.2e}")

# -------------------------------------------------------
# Hiện tượng tăng trưởng quá độ
# A = [[-1, 1000],[0, -1]], trị riêng đều = -1 (âm)
# nhưng ||e^At||_2 tăng mạnh trước khi giảm
# -------------------------------------------------------
A_nt = np.array([[-1.0, 1000.0],
                 [ 0.0,   -1.0]])

t_vals = np.linspace(0, 3, 300)
norms  = []
for t in t_vals:
    eAt = expm(A_nt * t)
    norms.append(np.linalg.norm(eAt, 2))

idx_max = np.argmax(norms)
print("\nHiện tượng tăng trưởng quá độ với A = [[-1,1000],[0,-1]]:")
print(f"  Giá trị lớn nhất của ||e^At||_2 = {norms[idx_max]:.1f}")
print(f"  Đạt tại t = {t_vals[idx_max]:.3f}")
print(f"  (dù rằng alpha(A) = -1 < 0 nên e^At phải giảm trong dài hạn)")

# Kiểm chứng bằng công thức giải tích
# e^{At} = e^{-t} [[1, 1000t],[0, 1]]
t_check = t_vals[idx_max]
eAt_analytic = np.exp(-t_check) * np.array([[1, 1000*t_check],
                                             [0, 1]])
print(f"  Kiểm tra: ||e^At||_2 (giải tích) = "
      f"{np.linalg.norm(eAt_analytic, 2):.1f}")
\end{lstlisting}

% ============================================================
% MỤC 4: HÀM DẤU, CĂN BẬC HAI VÀ LOGARIT MA TRẬN
% ============================================================
\section{Hàm dấu, căn bậc hai và logarit ma trận}
\label{sec:signlogsqrt}

\subsection{Hàm dấu ma trận (Matrix Sign Function)}

\subsubsection{Định nghĩa}

Với $z \in \C$ không nằm trên trục ảo:
\[
  \sign(z) = \begin{cases} -1 & \text{nếu } \re(z) < 0, \\ +1 & \text{nếu } \re(z) > 0. \end{cases}
\]

Giả sử $A \in \C^{n\times n}$ không có trị riêng thuần ảo, JCF được sắp xếp:
\[
  J = \begin{pmatrix} J_1 & 0 \\ 0 & J_2 \end{pmatrix},
\]
$\lambda(J_1)$ nằm trong nửa phẳng trái, $\lambda(J_2)$ nằm trong nửa phẳng phải. Thì:
\[
  \sign(A) = X \begin{pmatrix} -I_{m_1} & 0 \\ 0 & I_{m_2} \end{pmatrix} X^{-1}.
\]

\subsubsection{Ứng dụng tính không gian bất biến}

Với phân hoạch $X = [X_1\,|\,X_2]$, $X^{-H} = [Y_1\,|\,Y_2]$:
\[
  \frac{1}{2}(I + \sign(A)) = X_2 Y_2^H \quad \text{(rank $m_2$)}.
\]
Phân tích QR-với-đổi cột của ma trận này cho xấp xỉ không gian bất biến gắn với
các trị riêng nửa phẳng phải.

\subsubsection{Thuật toán lặp Newton}

Hàm dấu là nghiệm của $g(z) = z^2 - 1 = 0$. Tương tự Newton vô hướng:
\[
  S_0 = A; \qquad S_{k+1} = \frac{1}{2}\left(S_k + S_k^{-1}\right), \quad k = 0, 1, \ldots
\]

\begin{theorem}{Hội tụ Newton sign}{newtonsign}
Nếu $A$ không có trị riêng thuần ảo, dãy $\{S_k\}$ xác định với mọi $k$,
$\sign(S_k) = \sign(A)$ cho mọi $k$, và $S_k \to \sign(A)$ với tốc độ bậc hai:
\[
  \norm{S_{k+1} - S} \le \frac{1}{2}\norm{S_k^{-1}} \cdot \norm{S_k - S}^2.
\]
Bằng chứng dựa trên $G_k = (S_k - S)(S_k + S)^{-1}$ thỏa $G_{k+1} = G_k^2 \to 0$.
\end{theorem}

\subsubsection{Biến thể và gia tốc}

\textbf{Lặp Newton-Schulz:} Dùng xấp xỉ $S_k^{-1} \approx S_k(2I - S_k^2)$:
\[
  S_{k+1} = \frac{1}{2} S_k (3I - S_k^2).
\]
Không cần tính nghịch đảo ma trận, nhưng chỉ hội tụ cục bộ.

\textbf{Lặp có hệ số tỷ lệ:}
\[
  S_{k+1} = \frac{1}{2}\left(\mu_k S_k + (\mu_k S_k)^{-1}\right),
\]
với các lựa chọn $\mu_k = |\det(S_k)|^{1/n}$ hoặc $\mu_k = \sqrt{\rho(S_k^{-1})/\rho(S_k)}$.

\subsection{Căn bậc hai ma trận (Matrix Square Root)}

\subsubsection{Sự không đơn trị}

Ma trận có thể có nhiều căn bậc hai. Ví dụ:
\[
  A = \begin{pmatrix} 4 & 10 \\ 0 & 9 \end{pmatrix}
  \text{ có ít nhất 4 căn bậc hai: }
  \begin{pmatrix} \pm 2 & * \\ 0 & \pm 3 \end{pmatrix}.
\]

\begin{definition}{Căn bậc hai chính (Principal Square Root)}{principalsqrt}
$F = A^{1/2}$ là \emph{căn bậc hai chính} của $A$ nếu $F^2 = A$ và mọi trị riêng của
$F$ đều có phần thực dương.
\end{definition}

\subsubsection{Lặp Newton cho căn bậc hai}

Tương tự Newton vô hướng $x_{k+1} = (x_k + a/x_k)/2$:
\[
  X_0 = A; \qquad X_{k+1} = \frac{1}{2}(X_k + X_k^{-1}A), \quad k = 0, 1, \ldots
\]
Thay $X_k = A^{1/2} S_k$ vào trên cho ra lặp Newton sign cho $A^{1/2}$: hội tụ toàn
cục và bậc hai.

\subsubsection{Mối liên hệ với hàm dấu ma trận}

Áp dụng lặp Newton sign cho ma trận ghép:
\[
  \tilde{A} = \begin{pmatrix} 0 & A \\ I & 0 \end{pmatrix}.
\]
Qua quy nạp, các vòng lặp $\tilde{S}_k$ có dạng $\begin{pmatrix} 0 & X_k \\ Y_k & 0 \end{pmatrix}$,
và:
\[
  \sign\!\begin{pmatrix} 0 & A \\ I & 0 \end{pmatrix}
  = \begin{pmatrix} 0 & A^{1/2} \\ A^{-1/2} & 0 \end{pmatrix}.
\]

\subsubsection{Lặp Denman-Beavers}

Từ quan hệ trên, lặp Denman-Beavers có thuộc tính số tốt hơn Newton:
\[
  X_{k+1} = \frac{X_k + Y_k^{-1}}{2}, \quad
  Y_{k+1} = \frac{Y_k + X_k^{-1}}{2},
\]
với $X_0 = A$, $Y_0 = I$. Dãy này hội tụ $X_k \to A^{1/2}$, $Y_k \to A^{-1/2}$.

\subsection{Phân tích cực (Polar Decomposition)}

\subsubsection{Định lý cơ bản}

\begin{theorem}{Phân tích cực}{polar}
Nếu $A \in \R^{m \times n}$, $m \ge n$, thì tồn tại $U \in \R^{m \times n}$ (cột
trực chuẩn) và $P \in \R^{n \times n}$ (nửa xác định dương đối xứng) sao cho $A = UP$.
\end{theorem}

Chứng minh dựa trên SVD mỏng: $A = U_A \Sigma_A V_A^T$ cho $U = U_A V_A^T$,
$P = V_A \Sigma_A V_A^T$. Ta có $P = (A^T A)^{1/2}$ và nếu $\rank(A)=n$ thì
$U = A(A^T A)^{-1/2}$.

\subsubsection{Thuật toán lặp Newton cho nhân tử trực giao}

\[
  X_0 = A; \qquad X_{k+1} = \frac{X_k + X_k^{-T}}{2}, \quad k = 0, 1, \ldots
\]
Nếu $X_k = U_k P_k$ (phân tích cực của $X_k$), thì
\[
  X_{k+1} = U_k \cdot \frac{P_k + P_k^{-1}}{2}.
\]
Do $P_{k+1} = (P_k + P_k^{-1})/2$ đây chính là lặp Newton sign cho $P_0 = P$.
Vì $P_k \to \sign(P) = I$ bậc hai, nên $X_k \to U$ bậc hai.

\subsubsection{Kết nối với hàm dấu}

\[
  \sign\!\begin{pmatrix} 0 & A \\ A^T & 0 \end{pmatrix}
  = \begin{pmatrix} 0 & U \\ U^T & 0 \end{pmatrix},
\]
trong đó $U = U_A V_A^T$ là nhân tử trực giao của $A$.

\subsubsection{Lý thuyết nhiễu loạn}

\newpage
\begin{theorem}{Nhiễu loạn phân tích cực -- Li và Sun (2003)}{polarperturb}
Nếu $A, \tilde{A} \in \R^{n\times n}$ không suy biến với nhân tử trực giao $U$ và
$\tilde{U}$ tương ứng, thì
\[
  \norm{U - \tilde{U}}_F
  \le \frac{4\norm{A - \tilde{A}}_F}
           {\sigma_{n-1}(A) + \sigma_n(A) + \sigma_{n-1}(\tilde{A}) + \sigma_n(\tilde{A})}.
\]
\end{theorem}

\subsection{Logarit ma trận (Matrix Logarithm)}

\subsubsection{Định nghĩa và tính không đơn trị}

Nếu $e^X = A$, thì $X$ là một \emph{logarit} của $A$. Do tính tuần hoàn của hàm mũ
phức, $X + 2k\pi i$ cũng là logarit. Để khử mập mờ:

\begin{definition}{Logarit chính (Principal Logarithm)}{principallog}
Nếu $A \in \R^{n\times n}$ có mọi trị riêng thực đều dương, thì tồn tại duy nhất
ma trận thực $X$ thỏa $e^X = A$ với $\lambda(X) \subset \{z \in \C: -\pi < \im(z) < \pi\}$.
\end{definition}

\subsubsection{Các khai triển chuỗi cho hàm log}

\textbf{Chuỗi Maclaurin:}
\[
  \log(A) \approx M_q(A) = \sum_{k=1}^{q} \frac{(-1)^{k+1}(A-I)^k}{k},
  \quad \rho(A-I) < 1.
\]

\textbf{Chuỗi Gregory:}
\[
  \log(A) \approx G_q(A) = -2\sum_{k=0}^{q} \frac{1}{2k+1}
  \left[(I-A)(I+A)^{-1}\right]^{2k+1},
  \quad \re(\lambda_i) > 0.
\]

\textbf{Xấp xỉ Padé chéo:} Ví dụ xấp xỉ $(3,3)$:
\[
  \log(A) \approx r_{33}(A) = D(A)^{-1} N(A),
\]
\[
  D(A) = 60I + 90(A-I) + 36(A-I)^2 + 3(A-I)^3,
\]
\[
  N(A) = 60(A-I) + 60(A-I)^2 + 11(A-I)^3.
\]

\subsubsection{Thuật toán Inverse Scaling and Squaring}

\begin{algobox}
\textbf{Thuật toán 1.4.1 (Inverse Scaling and Squaring cho $\log(A)$)}

\medskip
1. Lấy căn bậc hai lặp cho đến $\norm{A_k - I} \le \text{tol}$:
\[
   k = 0,\quad A_0 = A;\quad
   \text{Khi } \norm{A-I} > \text{tol}: k \leftarrow k+1,\quad A_k = A_{k-1}^{1/2}.
\]
(Dùng lặp Denman-Beavers để tính căn bậc hai.)

2. Tính $F \approx \log(A_k)$ bằng xấp xỉ Padé phù hợp.

3. Hoàn nguyên: $\log(A) = 2^k F$.

\medskip
Đây là quá trình ngược của Scaling and Squaring cho hàm mũ.
\end{algobox}

% ------------------------------------------------------------------
% PYTHON MINH HỌA – HÀM DẤU, CĂN BẬC HAI, LOGARIT
% ------------------------------------------------------------------
\subsection{Mã nguồn Python minh họa}

Mã nguồn sau minh họa bốn thuật toán: (1) lặp Newton cho hàm dấu ma trận,
(2) lặp Denman-Beavers cho căn bậc hai, (3) lặp Newton cho nhân tử trực giao
(phân tích cực), và (4) thuật toán Inverse Scaling and Squaring cho logarit.

\begin{lstlisting}[style=python, caption={Hàm dấu, căn bậc hai, phân tích cực và logarit ma trận -- lặp Newton và Denman-Beavers}]
import numpy as np
from scipy.linalg import sqrtm, logm, svd

# -------------------------------------------------------
# 1. Lặp Newton cho hàm dấu ma trận
#    S_{k+1} = 0.5*(S_k + S_k^{-1})
# -------------------------------------------------------
def matrix_sign_newton(A, max_iter=50, tol=1e-14):
    S = A.astype(complex).copy()
    for k in range(max_iter):
        S_inv = np.linalg.inv(S)
        S_new = 0.5 * (S + S_inv)
        residual = np.linalg.norm(S_new - S, 'fro')
        if residual < tol:
            print(f"  [sign] Hội tụ sau {k+1} vòng lặp, phần dư = {residual:.2e}")
            return S_new
        S = S_new
    print(f"  [sign] Chưa hội tụ sau {max_iter} vòng lặp")
    return S

# -------------------------------------------------------
# 2. Lặp Denman-Beavers cho căn bậc hai ma trận
#    X_{k+1} = 0.5*(X_k + Y_k^{-1})
#    Y_{k+1} = 0.5*(Y_k + X_k^{-1})
# -------------------------------------------------------
def matrix_sqrt_denman_beavers(A, max_iter=50, tol=1e-14):
    X = A.astype(float).copy()
    Y = np.eye(A.shape[0])
    for k in range(max_iter):
        X_inv = np.linalg.inv(X)
        Y_inv = np.linalg.inv(Y)
        X_new = 0.5 * (X + Y_inv)
        Y_new = 0.5 * (Y + X_inv)
        err = np.linalg.norm(X_new @ X_new - A, 'fro')
        if err < tol:
            print(f"  [sqrt] Hội tụ sau {k+1} vòng lặp, ||X^2-A||_F = {err:.2e}")
            return X_new, Y_new
        X, Y = X_new, Y_new
    return X, Y

# -------------------------------------------------------
# 3. Phân tích cực: Newton lặp cho nhân tử trực giao
#    X_{k+1} = 0.5*(X_k + X_k^{-T})
# -------------------------------------------------------
def polar_newton(A, max_iter=50, tol=1e-14):
    X = A.astype(float).copy()
    for k in range(max_iter):
        X_invT = np.linalg.inv(X).T
        X_new = 0.5 * (X + X_invT)
        err = np.linalg.norm(X_new.T @ X_new - np.eye(A.shape[0]), 'fro')
        if err < tol:
            print(f"  [polar] Hội tụ sau {k+1} vòng lặp, ||U^T U - I||_F = {err:.2e}")
            return X_new
        X = X_new
    return X

# -------------------------------------------------------
# 4. Logarit ma trận: Inverse Scaling and Squaring
#    Dùng căn bậc hai lặp + Maclaurin cho log
# -------------------------------------------------------
def matrix_log_iss(A, tol_norm=0.5, q=20):
    """Inverse Scaling and Squaring cho logm."""
    Ak = A.astype(float).copy()
    k  = 0
    while np.linalg.norm(Ak - np.eye(Ak.shape[0]), 'fro') > tol_norm:
        Ak, _ = matrix_sqrt_denman_beavers(Ak, tol=1e-13)
        k += 1
    # Tính log(Ak) bằng chuỗi Maclaurin (Ak gần I)
    B   = Ak - np.eye(Ak.shape[0])
    logA = np.zeros_like(B)
    Bk   = np.eye(Ak.shape[0])
    for j in range(1, q + 1):
        Bk    = Bk @ B
        logA += ((-1)**(j+1) / j) * Bk
    return (2**k) * logA   # hoàn nguyên

# -------------------------------------------------------
# ỨNG DỤNG VÀ KIỂM TRA
# -------------------------------------------------------
A = np.array([[4.0, 3.0],
              [1.0, 2.0]])

print("Ma trận A:")
print(A)

# Hàm dấu
S = matrix_sign_newton(A)
print("\nHàm dấu sign(A):")
print(S.real)

# Căn bậc hai
X_sqrt, _ = matrix_sqrt_denman_beavers(A)
ref_sqrt   = sqrtm(A)
print("\nCăn bậc hai A^{1/2} (Denman-Beavers):")
print(X_sqrt)
print(f"  Sai số so với sqrtm(scipy): {np.linalg.norm(X_sqrt - ref_sqrt, 'fro'):.2e}")

# Phân tích cực
U_polar = polar_newton(A)
U_svd   = (lambda U,_,Vt: U @ Vt)(*svd(A, full_matrices=False))
print("\nNhân tử trực giao U (phân tích cực, Newton):")
print(U_polar)
print(f"  Sai số so với SVD: {np.linalg.norm(U_polar - U_svd, 'fro'):.2e}")

# Logarit
L = matrix_log_iss(A, tol_norm=0.5, q=30)
ref_log = logm(A)
print("\nLogarit chính log(A) (Inverse Scaling and Squaring):")
print(L)
print(f"  Sai số so với logm(scipy): {np.linalg.norm(L - ref_log, 'fro'):.2e}")
print(f"  Kiểm tra: ||e^L - A||_F = "
      f"{np.linalg.norm(np.linalg.matrix_power(np.eye(2), 0) + L - A, 'fro'):.2e}")
\end{lstlisting}

% ============================================================
% CHƯƠNG 2: TRIỂN KHAI THỰC NGHIỆM
% ============================================================
\chapter{Triển khai thực nghiệm}

Chương này trình bày các thực nghiệm tính toán cụ thể, đánh giá định lượng hiệu năng
và độ chính xác của các phương pháp đã trình bày ở Chương~1. Tất cả thực nghiệm được
thực hiện bằng Python~3.10 với các thư viện NumPy~1.25 và SciPy~1.11 trên môi trường
chuẩn IEEE 754 double precision ($u \approx 2.22 \times 10^{-16}$).

\section{Môi trường và thiết lập thực nghiệm}

\subsection{Yêu cầu thư viện}

Toàn bộ mã thực nghiệm sử dụng các thư viện Python phổ biến, dễ dàng cài đặt:

\begin{lstlisting}[style=python, caption={Thiết lập môi trường thực nghiệm}]
# Cài đặt thư viện (nếu chưa có)
# pip install numpy scipy matplotlib

import numpy as np
from scipy.linalg import (expm, logm, sqrtm,
                           schur, svd, norm)
import time

# Thiết lập random seed để tái hiện kết quả
np.random.seed(42)

# Epsilon máy
u = np.finfo(float).eps
print(f"Epsilon máy u = {u:.3e}")
print(f"Phiên bản NumPy: {np.__version__}")

# Hàm tạo ma trận ngẫu nhiên với số điều kiện cho trước
def random_matrix_cond(n, kappa):
    """Tạo ma trận ngẫu nhiên n x n với condition number xấp xỉ kappa."""
    U, _ = np.linalg.qr(np.random.randn(n, n))
    V, _ = np.linalg.qr(np.random.randn(n, n))
    s = np.logspace(0, -np.log10(kappa), n)
    return U @ np.diag(s) @ V.T

# Hàm đo sai số tương đối Frobenius
def rel_err(F_approx, F_exact):
    denom = norm(F_exact, 'fro')
    return norm(F_approx - F_exact, 'fro') / denom if denom > 0 else float('inf')
\end{lstlisting}

\subsection{Các ma trận thử nghiệm chuẩn}

Để đánh giá toàn diện, chúng ta dùng bốn lớp ma trận thử nghiệm đại diện cho
các trường hợp khó khác nhau trong thực tế.

\begin{lstlisting}[style=python, caption={Định nghĩa các ma trận thử nghiệm}]
def test_matrices():
    """Trả về dict các ma trận thử nghiệm chuẩn."""
    mats = {}

    # 1. Ma trận chéo hóa được, điều kiện tốt
    mats['diagonal'] = np.diag([1.0, 2.0, 3.0, 4.0])

    # 2. Ma trận đối xứng xác định dương (SPD)
    B = np.random.randn(4, 4)
    mats['SPD'] = B.T @ B + 4 * np.eye(4)

    # 3. Ma trận có trị riêng phức (xoay)
    theta = np.pi / 4
    mats['rotation'] = np.array([
        [ np.cos(theta), -np.sin(theta), 0, 0],
        [ np.sin(theta),  np.cos(theta), 0, 0],
        [0, 0,  np.cos(theta), -np.sin(theta)],
        [0, 0,  np.sin(theta),  np.cos(theta)]
    ])

    # 4. Ma trận không chính tắc (non-normal) - khá khó
    mats['non_normal'] = np.array([
        [-1.0, 500.0, 0.0, 0.0],
        [ 0.0,  -1.0, 500.0, 0.0],
        [ 0.0,   0.0, -1.0, 500.0],
        [ 0.0,   0.0,  0.0, -1.0]
    ])

    # 5. Ma trận trị riêng lặp (Jordan block 4x4)
    lam = 2.0
    mats['jordan'] = lam * np.eye(4) + np.diag([1, 1, 1], k=1)

    return mats

MATS = test_matrices()
for name, M in MATS.items():
    kappa = np.linalg.cond(M)
    eigvals = np.linalg.eigvals(M)
    print(f"  {name:15s}: kích thước={M.shape}, "
          f"kappa={kappa:.2e}, "
          f"trị riêng ~ {np.round(eigvals, 2)}")
\end{lstlisting}

\section{Thực nghiệm 1: Phương pháp trị riêng và Schur-Parlett}

Thực nghiệm này so sánh ba phương pháp tính $e^A$: phân tích trị riêng (eigenvector),
Schur-Parlett, và SciPy chuẩn. Chúng ta đo sai số tương đối Frobenius và thời gian chạy.

\begin{lstlisting}[style=python, caption={Thực nghiệm 1: so sánh phương pháp trị riêng vs Schur-Parlett}]
import numpy as np
from scipy.linalg import expm, schur
import time

def expm_eigenvector(A):
    """Tính e^A qua phân tích trị riêng."""
    lam, X = np.linalg.eig(A)
    return X @ np.diag(np.exp(lam)) @ np.linalg.inv(X)

def expm_schur_parlett(A):
    """
    Tính e^A qua Schur-Parlett (đơn giản hóa cho trị riêng phân biệt).
    A = Q T Q^H  =>  e^A = Q e^T Q^H
    """
    T, Q = schur(A, output='complex')
    n    = T.shape[0]
    F    = np.zeros_like(T, dtype=complex)
    lam  = np.diag(T)

    # Khởi tạo đường chéo
    for i in range(n):
        F[i, i] = np.exp(lam[i])

    # Tính các siêu đường chéo theo thứ tự p = 1, 2, ..., n-1
    for p in range(1, n):
        for i in range(n - p):
            j = i + p
            dl = lam[j] - lam[i]
            if abs(dl) > 1e-10:
                # Sai phân chia bậc 1
                Fij = T[i, j] * (F[j, j] - F[i, i]) / dl
            else:
                # Giới hạn: f'(lambda) khi trị riêng trùng nhau
                Fij = T[i, j] * np.exp(lam[i])
            # Đóng góp từ các phần tử trung gian
            for k in range(i + 1, j):
                dlk = lam[j] - lam[k]
                dki = lam[k] - lam[i]
                if abs(dlk) > 1e-10 and abs(dki) > 1e-10:
                    Fij += (T[i, k] * F[k, j] - F[i, k] * T[k, j]) / dl
            F[i, j] = Fij

    return (Q @ F @ Q.conj().T).real

# -------------------------------------------------------
# Chạy thử nghiệm trên các ma trận chuẩn
# -------------------------------------------------------
print(f"{'Ma trận':15s} | {'Sai số Eig':10s} | {'Sai số Schur':10s} "
      f"| {'t_eig (ms)':11s} | {'t_schur (ms)':12s}")
print("-" * 68)

for name, A in MATS.items():
    ref = expm(A)

    t0 = time.perf_counter()
    F_eig = expm_eigenvector(A).real
    t_eig = (time.perf_counter() - t0) * 1000

    t0 = time.perf_counter()
    F_sch = expm_schur_parlett(A)
    t_sch = (time.perf_counter() - t0) * 1000

    e_eig = rel_err(F_eig, ref)
    e_sch = rel_err(F_sch, ref)

    print(f"{name:15s} | {e_eig:10.2e} | {e_sch:10.2e} "
          f"| {t_eig:11.3f} | {t_sch:12.3f}")
\end{lstlisting}

\section{Thực nghiệm 2: Phương pháp xấp xỉ -- Taylor và ảnh hưởng của hủy số}

Thực nghiệm này khảo sát kỹ lưỡng hiện tượng catastrophic cancellation trong
xấp xỉ Taylor, và minh họa tại sao bậc xấp xỉ cao hơn không luôn tốt hơn.

\begin{lstlisting}[style=python, caption={Thực nghiệm 2: hiện tượng hủy số trong Taylor và vai trò của scaling}]
import numpy as np
from scipy.linalg import expm

def expm_taylor(A, q):
    n = A.shape[0]
    F = np.eye(n); Ak = np.eye(n); fact = 1.0
    for k in range(1, q + 1):
        Ak = Ak @ A; fact *= k; F += Ak / fact
    return F

def expm_taylor_scaled(A, q, j=None):
    """Taylor với scaling: tính e^{A/2^j} rồi bình phương j lần."""
    norm_A = np.max(np.sum(np.abs(A), axis=1))
    if j is None:
        j = max(0, int(np.ceil(np.log2(norm_A))))
    As = A / (2 ** j)
    F = expm_taylor(As, q)
    for _ in range(j): F = F @ F
    return F

# Ma trận gây hủy số nhất
A_bad = np.array([[-49.0, 24.0],
                  [-64.0, 31.0]])
ref = expm(A_bad)

print("Hiện tượng hủy số với A = [[-49,24],[-64,31]]")
print(f"  Kết quả chính xác: \n{ref}\n")

print(f"  {'Bậc q':6s} | {'Taylor thuần (err)':20s} | {'Taylor + Scaling (err)':22s}")
print("  " + "-"*55)
for q in [5, 10, 15, 20, 30, 50, 59]:
    F_naive  = expm_taylor(A_bad, q)
    F_scaled = expm_taylor_scaled(A_bad, q)
    e_naive  = rel_err(F_naive,  ref)
    e_scaled = rel_err(F_scaled, ref)
    print(f"  {q:6d} | {e_naive:20.4e} | {e_scaled:22.4e}")

# Theo dõi giá trị lớn nhất của các số hạng trung gian
print("\nGiá trị lớn nhất của các số hạng Ak/k! (A_bad):")
n = A_bad.shape[0]
Ak = np.eye(n); fact = 1.0
for k in range(1, 21):
    Ak = Ak @ A_bad; fact *= k
    term_max = np.max(np.abs(Ak / fact))
    if k <= 5 or k % 5 == 0:
        print(f"  k={k:2d}: max|A^k / k!| = {term_max:.3e}")
\end{lstlisting}

\section{Thực nghiệm 3: Scaling and Squaring -- phân tích hiệu năng theo kích thước}

Thực nghiệm này đo thời gian chạy và độ chính xác của thuật toán Scaling and Squaring
(SciPy expm) so với phương pháp trị riêng thuần túy khi kích thước ma trận $n$ tăng.

\begin{lstlisting}[style=python, caption={Thực nghiệm 3: hiệu năng Scaling and Squaring theo kích thước n}]
import numpy as np
from scipy.linalg import expm
import time

def expm_eigenvector_general(A):
    lam, X = np.linalg.eig(A)
    return (X @ np.diag(np.exp(lam)) @ np.linalg.inv(X)).real

sizes = [4, 8, 16, 32, 64, 128, 256]
n_repeat = 5   # số lần lặp để lấy trung bình

print(f"{'n':>6} | {'t_eig (ms)':>12} | {'t_expm (ms)':>12} "
      f"| {'Tăng tốc':>8} | {'sai số eig':>10}")
print("-" * 60)

for n in sizes:
    # Tạo ma trận ngẫu nhiên xác định dương (luôn có e^A xác định)
    B = np.random.randn(n, n)
    A = (B + B.T) / 2 - n * np.eye(n)   # đối xứng, âm định

    ref = expm(A)

    # Đo thời gian expm (SciPy -- Padé + Scaling&Squaring)
    t_expm_list = []
    for _ in range(n_repeat):
        t0 = time.perf_counter()
        F_expm = expm(A)
        t_expm_list.append((time.perf_counter() - t0) * 1000)
    t_expm = np.median(t_expm_list)

    # Đo thời gian phương pháp trị riêng (chỉ khả thi khi n nhỏ)
    if n <= 64:
        t_eig_list = []
        for _ in range(n_repeat):
            t0 = time.perf_counter()
            F_eig = expm_eigenvector_general(A)
            t_eig_list.append((time.perf_counter() - t0) * 1000)
        t_eig   = np.median(t_eig_list)
        err_eig = rel_err(F_eig, ref)
        speedup = t_eig / t_expm
        print(f"{n:>6} | {t_eig:>12.3f} | {t_expm:>12.3f} "
              f"| {speedup:>8.2f}x | {err_eig:>10.2e}")
    else:
        print(f"{n:>6} | {'N/A (quá lớn)':>12} | {t_expm:>12.3f} "
              f"| {'---':>8} | {'---':>10}")
\end{lstlisting}

\section{Thực nghiệm 4: Hàm dấu, căn bậc hai và logarit -- hội tụ và độ chính xác}

Thực nghiệm cuối cùng phân tích chi tiết tốc độ hội tụ của lặp Newton và
Denman-Beavers, đồng thời so sánh kết quả với hàm chuẩn SciPy.

\begin{lstlisting}[style=python, caption={Thực nghiệm 4: hội tụ bậc hai của lặp Newton và Denman-Beavers}]
import numpy as np
from scipy.linalg import sqrtm, logm, expm

# Cài đặt lại các hàm lặp với theo dõi hội tụ
def matrix_sign_newton_track(A, max_iter=30, tol=1e-14):
    S = A.astype(complex).copy()
    residuals = []
    for k in range(max_iter):
        S_new = 0.5 * (S + np.linalg.inv(S))
        res   = np.linalg.norm(S_new @ S_new - np.eye(S.shape[0]), 'fro')
        residuals.append(res)
        if res < tol: S = S_new; break
        S = S_new
    return S, residuals

def matrix_sqrt_db_track(A, max_iter=30, tol=1e-14):
    X = A.astype(float).copy()
    Y = np.eye(A.shape[0])
    residuals = []
    for k in range(max_iter):
        X_inv = np.linalg.inv(X)
        Y_inv = np.linalg.inv(Y)
        X_new = 0.5 * (X + Y_inv)
        Y_new = 0.5 * (Y + X_inv)
        res   = np.linalg.norm(X_new @ X_new - A, 'fro')
        residuals.append(res)
        if res < tol: X, Y = X_new, Y_new; break
        X, Y = X_new, Y_new
    return X, Y, residuals

# -------------------------------------------------------
# Ma trận thử nghiệm: SPD, để kiểm chứng
# -------------------------------------------------------
B = np.array([[4.0, 2.0, 1.0],
              [2.0, 5.0, 3.0],
              [1.0, 3.0, 6.0]])

print("=" * 55)
print("MA TRẬN B =")
print(B)
print(f"Trị riêng: {np.linalg.eigvals(B).real}")

# Thí nghiệm: Hàm dấu sign(B)
print("\n--- Hàm dấu ma trận (Newton) ---")
S_B, res_sign = matrix_sign_newton_track(B)
ref_sign = np.diag(np.sign(np.linalg.eigvals(B).real))
# Lưu ý: sign(B) phụ thuộc vào trị riêng thực (đều dương => sign = I)
print(f"  sign(B) (trị riêng đều dương, nên sign(B) = I):")
print(f"  {np.diag(S_B.real)}")
print("  Phần dư ||S^2 - I||_F theo vòng lặp:")
for k, r in enumerate(res_sign):
    print(f"    Vòng {k+1}: {r:.4e}")

# Thí nghiệm: Căn bậc hai Denman-Beavers
print("\n--- Căn bậc hai ma trận (Denman-Beavers) ---")
X_sqrt, _, res_sqrt = matrix_sqrt_db_track(B)
ref_sqrt = sqrtm(B)
err_sqrt = np.linalg.norm(X_sqrt - ref_sqrt, 'fro') / np.linalg.norm(ref_sqrt, 'fro')
print(f"  Sai số so với sqrtm(scipy): {err_sqrt:.2e}")
print("  Phần dư ||X^2 - B||_F theo vòng lặp:")
for k, r in enumerate(res_sqrt):
    print(f"    Vòng {k+1}: {r:.4e}")

# Thí nghiệm: Logarit -- kiểm chứng e^{logm(B)} = B
print("\n--- Logarit ma trận (SciPy logm) ---")
L_B   = logm(B)
check = expm(L_B)
err_log = np.linalg.norm(check - B, 'fro') / np.linalg.norm(B, 'fro')
print(f"  Sai số ||e^{{logm(B)}} - B||_F / ||B||_F = {err_log:.2e}")

# -------------------------------------------------------
# Thí nghiệm phân tích cực với SVD đối chiếu
# -------------------------------------------------------
print("\n--- Phân tích cực A = U P ---")
A_rect = np.array([[1.0, 2.0, 0.0],
                   [0.0, 3.0, 4.0],
                   [5.0, 0.0, 1.0]])
Ua, Sa, Vat = np.linalg.svd(A_rect)
U_polar_ref = Ua @ Vat          # nhân tử trực giao từ SVD
P_polar_ref = Vat.T @ np.diag(Sa) @ Vat   # nhân tử xác định dương

# Newton lặp cho nhân tử trực giao
X = A_rect.copy()
for k in range(30):
    X_new = 0.5 * (X + np.linalg.inv(X).T)
    if np.linalg.norm(X_new - X, 'fro') < 1e-14: break
    X = X_new
U_newton = X

err_polar = np.linalg.norm(U_newton - U_polar_ref, 'fro')
print(f"  ||U_Newton - U_SVD||_F = {err_polar:.2e}")
print(f"  Kiểm tra trực giao: ||U^T U - I||_F = "
      f"{np.linalg.norm(U_newton.T @ U_newton - np.eye(3), 'fro'):.2e}")
print(f"  Kiểm tra phân tích: ||A - UP||_F = "
      f"{np.linalg.norm(A_rect - U_newton @ P_polar_ref, 'fro'):.2e}")
\end{lstlisting}

\section{Tổng hợp và nhận xét kết quả}

Dựa trên bốn thực nghiệm trên, ta có thể đưa ra một số nhận xét tổng hợp.

\subsection{So sánh tổng quan các phương pháp}

\begin{center}
\resizebox{\textwidth}{!}{%
\begin{tabular}{lcccc}
\toprule
\textbf{Phương pháp} & \textbf{Độ chính xác} & \textbf{Ổn định số} & \textbf{Chi phí} & \textbf{Phù hợp khi} \\
\midrule
Trị riêng (eigenvector) & Phụ thuộc $\kappa_2(X)$ & Kém nếu $X$ ill-cond & $O(n^3)$ & $A$ chính tắc \\
Schur-Parlett           & Cao                    & Tốt                  & $O(n^3)$ & $n$ nhỏ-vừa \\
Taylor cắt cụt          & Bị hủy số              & Kém                  & $O(qn^3)$ & Chỉ để minh họa \\
Taylor + Scaling        & Tốt                    & Khá                  & $O(qn^3)$ & $\norm{A}$ lớn \\
Padé + Scal. \& Sq.    & Cao                    & Rất tốt              & $O(n^3)$ & Mục đích chung \\
Newton sign             & Bậc hai                & Tốt                  & $O(kn^3)$ & Hàm dấu, căn \\
Denman-Beavers          & Bậc hai                & Tốt hơn Newton       & $O(kn^3)$ & Căn bậc hai \\
\bottomrule
\end{tabular}}
\end{center}

\subsection{Khuyến nghị thực tế}

\begin{importantbox}
\textbf{Khuyến nghị thực hành:}
\begin{enumerate}
  \item \textbf{Hàm mũ $e^A$:} Luôn dùng \texttt{scipy.linalg.expm}
    (Padé bậc 13 + Scaling and Squaring) trong thực tế. Chỉ tự cài đặt
    khi nghiên cứu hoặc yêu cầu đặc biệt về hiệu năng.

  \item \textbf{Logarit $\log(A)$:} Dùng \texttt{scipy.linalg.logm}
    (Inverse Scaling and Squaring). Cẩn thận khi $A$ có trị riêng âm hoặc
    phức (vấn đề chọn nhánh logarit).

  \item \textbf{Căn bậc hai $A^{1/2}$:} \texttt{scipy.linalg.sqrtm} cho kết quả
    tốt. Lặp Denman-Beavers hữu ích khi cần kiểm soát trực tiếp quá trình hội tụ.

  \item \textbf{Ma trận lớn ($n \gg 1000$):} Chuyển sang phương pháp Krylov
    (ví dụ: Expokit, Expmv) thay vì các phương pháp dense.

  \item \textbf{Kiểm tra điều kiện:} Trước khi tính bất kỳ hàm ma trận nào,
    nên kiểm tra số điều kiện $\kappa_2(A)$ và hoành độ phổ $\alpha(A)$ để dự
    báo độ chính xác kết quả.
\end{enumerate}
\end{importantbox}

\subsection{Tổng kết}

Qua lý thuyết ở Chương~1 và thực nghiệm ở Chương~2, ta thấy rõ rằng không có
một phương pháp duy nhất nào tối ưu cho mọi trường hợp. Sự lựa chọn phụ thuộc
vào cấu trúc ma trận (chính tắc hay không, điều kiện tốt hay xấu, kích thước),
hàm cần tính ($e^A$, $\log A$, $A^{1/2}$, $\sign(A)$), và yêu cầu về độ chính xác
và tốc độ. Tuy nhiên, các phương pháp dựa trên phân tích Schur và xấp xỉ Padé
với Scaling and Squaring nổi bật là lựa chọn ổn định và hiệu quả nhất cho mục
đích tổng quát, và đây cũng là nền tảng của các thư viện số hiện đại như MATLAB,
SciPy và Julia.

\end{document}
